\documentclass[10pt,a4paper]{amsart}
\usepackage{amsmath,amssymb,amsthm,graphicx}
\usepackage{verbatim}
\newenvironment{reference}{\comment}{\endcomment}
%\newenvironment{reference}{}{}
\newenvironment{slogan}{\comment}{\endcomment}
\newenvironment{history}{\comment}{\endcomment}
% For commutative diagrams we use Xy-pic
\usepackage[all]{xy}

% We use 2cell for 2-commutative diagrams.
\xyoption{2cell}
% Theorem environments.
%
\theoremstyle{plain}
\newtheorem{theorem}[subsection]{Theorem}
\newtheorem{proposition}[subsection]{Proposition}
\newtheorem{lemma}[subsection]{Lemma}

\theoremstyle{definition}
\newtheorem{definition}[subsection]{Definition}
\newtheorem{example}[subsection]{Example}
\newtheorem{exercise}[subsection]{Exercise}
\newtheorem{situation}[subsection]{Situation}

\theoremstyle{remark}
\newtheorem{remark}[subsection]{Remark}
\newtheorem{remarks}[subsection]{Remarks}

\numberwithin{equation}{subsection}

% Macros
%
\def\lim{\mathop{\mathrm{lim}}\nolimits}
\def\colim{\mathop{\mathrm{colim}}\nolimits}
\def\Spec{\mathop{\mathrm{Spec}}}
\def\Hom{\mathop{\mathrm{Hom}}\nolimits}
\def\Ext{\mathop{\mathrm{Ext}}\nolimits}
\def\SheafHom{\mathop{\mathcal{H}\!\mathit{om}}\nolimits}
\def\SheafExt{\mathop{\mathcal{E}\!\mathit{xt}}\nolimits}
\def\Sch{\mathit{Sch}}
\def\Mor{\mathop{\mathrm{Mor}}\nolimits}
\def\Ob{\mathop{\mathrm{Ob}}\nolimits}
\def\Sh{\mathop{\mathit{Sh}}\nolimits}
\def\NL{\mathop{N\!L}\nolimits}
\def\CH{\mathop{\mathrm{CH}}\nolimits}
\def\proetale{{pro\text{-}\acute{e}tale}}
\def\etale{{\acute{e}tale}}
\def\QCoh{\mathit{QCoh}}
\def\Ker{\mathop{\mathrm{Ker}}}
\def\Im{\mathop{\mathrm{Im}}}
\def\Coker{\mathop{\mathrm{Coker}}}
\def\Coim{\mathop{\mathrm{Coim}}}

% Boxtimes
%
\DeclareMathSymbol{\boxtimes}{\mathbin}{AMSa}{"02}

%
% Macros for moduli stacks/spaces
%
\def\QCohstack{\mathcal{QC}\!\mathit{oh}}
\def\Cohstack{\mathcal{C}\!\mathit{oh}}
\def\Spacesstack{\mathcal{S}\!\mathit{paces}}
\def\Quotfunctor{\mathrm{Quot}}
\def\Hilbfunctor{\mathrm{Hilb}}
\def\Curvesstack{\mathcal{C}\!\mathit{urves}}
\def\Polarizedstack{\mathcal{P}\!\mathit{olarized}}
\def\Complexesstack{\mathcal{C}\!\mathit{omplexes}}
% \Pic is the operator that assigns to X its picard group, usage \Pic(X)
% \Picardstack_{X/B} denotes the Picard stack of X over B
% \Picardfunctor_{X/B} denotes the Picard functor of X over B
\def\Pic{\mathop{\mathrm{Pic}}\nolimits}
\def\Picardstack{\mathcal{P}\!\mathit{ic}}
\def\Picardfunctor{\mathrm{Pic}}
\def\Deformationcategory{\mathcal{D}\!\mathit{ef}}
\usepackage{cleveref}
\crefname{originallemma}{original lemma}{original lemmas}
\Crefname{originallemma}{Original lemma}{Original lemmas}
\makeatletter
\newlabel{coherent-lemma-quasi-coherent-affine-cohomology-zero}{{01XB}{source}{}{sourcepointer-1}{}}
\newlabel{coherent-lemma-quasi-coherent-affine-cohomology-zero@cref}{{[originallemma][0][]01XB}{[1][0][]source}{}{}{}}
\makeatother
\crefname{originalproposition}{original proposition}{original propositions}
\Crefname{originalproposition}{Original proposition}{Original propositions}
\makeatletter
\newlabel{groupoids-proposition-finite-flat-equivalence}{{03BM}{source}{}{sourcepointer-2}{}}
\newlabel{groupoids-proposition-finite-flat-equivalence@cref}{{[originalproposition][0][]03BM}{[1][0][]source}{}{}{}}
\makeatother
\crefname{originallemma}{original lemma}{original lemmas}
\Crefname{originallemma}{Original lemma}{Original lemmas}
\makeatletter
\newlabel{lemma-alternating-spectral-sequence}{{0728}{source}{}{sourcepointer-3}{}}
\newlabel{lemma-alternating-spectral-sequence@cref}{{[originallemma][0][]0728}{[1][0][]source}{}{}{}}
\makeatother
\crefname{originallemma}{original lemma}{original lemmas}
\Crefname{originallemma}{Original lemma}{Original lemmas}
\makeatletter
\newlabel{lemma-quasi-coherent-twist}{{072A}{source}{}{sourcepointer-4}{}}
\newlabel{lemma-quasi-coherent-twist@cref}{{[originallemma][0][]072A}{[1][0][]source}{}{}{}}
\makeatother
\crefname{originalsection}{original section}{original sections}
\Crefname{originalsection}{Original section}{Original sections}
\makeatletter
\newlabel{section-higher-direct-image}{{071Y}{source}{}{sourcepointer-5}{}}
\newlabel{section-higher-direct-image@cref}{{[originalsection][0][]071Y}{[1][0][]source}{}{}{}}
\makeatother
\crefname{originallemma}{original lemma}{original lemmas}
\Crefname{originallemma}{Original lemma}{Original lemmas}
\makeatletter
\newlabel{spaces-morphisms-lemma-compose-after-separated}{{03KR}{source}{}{sourcepointer-6}{}}
\newlabel{spaces-morphisms-lemma-compose-after-separated@cref}{{[originallemma][0][]03KR}{[1][0][]source}{}{}{}}
\makeatother
\title{Vanishing above the number of sheets of an affine etale cover of a quasi-compact separated space}
\author{The Stacks Project Authors}
\date{stacks-project:spaces-cohomology.tex:local-snapshot-20261012; GNU FDL 1.2 or later}
\begin{document}
\maketitle
\noindent\textit{Editorial note.} All delivered mathematical text is the author's own, copied from the pinned Stacks Project chapter file spaces-cohomology.tex; no statement or proof step was rewritten, reordered or removed.\par\medskip
\noindent\textit{Editorial note.} Preamble provenance: the pinned chapter file opens with \textbackslash{}input\{preamble\} and carries no preamble text of its own, so the project's single Stacks-wide preamble staged at sources/stacks-spaces-cohomology/preamble.tex is excerpted verbatim as three spans, lines 8-12 (verbatim package and the project's own reference/slogan/history display environments), lines 14-18 (Xy-pic setup) and lines 154-217 (theorem declarations and mathematical macros). The wrapper class and the amsmath/amssymb/amsthm/graphicx base packages are supplied by the assembler. No macro was invented or redefined.\par\medskip
\noindent\textit{Editorial note.} Same-chapter pointers are located at their \textbackslash{}label lines inside the pinned chapter file; cross-chapter pointers are located in the project's own tag list staged in the snapshot (sources/stacks-chow/tags.txt), at the verbatim lines carrying their canonical tags.\par\medskip
\noindent\textit{Allowed standard prerequisites (editorial).}\begin{itemize}
\item The alternating Cech spectral sequence attached to a surjective etale morphism. Same chapter spaces-cohomology.tex, label lemma-alternating-spectral-sequence (invoked at proof line 1345). It computes the cohomology of \$X\$ from the twisted cohomology of the configuration quotients \$U\_p\$, and requires \$f\$ to be separated, which holds because \$U\$ and \$X\$ are separated.
\item Separatedness of the composite, and the finite flat equivalence for groupoid quotients. Morphisms of Spaces, label spaces-morphisms-lemma-compose-after-separated (invoked at proof line 1348) and Groupoids, label groupoids-proposition-finite-flat-equivalence (invoked at proof line 1354). The first gives that \$f\$ is separated; the second identifies \$U\_p = W\_p/S\_\{p+1\}\$ as an affine scheme because \$W\_p\$ is affine and the action is free and finite.
\item Vanishing of the higher cohomology of a quasi-coherent module on the affine configuration scheme, and quasi-coherence of the twists. Cohomology of Schemes, label coherent-lemma-quasi-coherent-affine-cohomology-zero (invoked at proof line 1361) and Same chapter spaces-cohomology.tex, label lemma-quasi-coherent-twist (invoked at proof line 1363). They make the \$E\_1\$-terms vanish in positive degrees: the twisted sheaves are quasi-coherent, their cohomology on \$W\_p\$ agrees with the cohomology of the corresponding module on the affine scheme, and that vanishes above degree zero.
\end{itemize}\medskip
\ifdefined\hypertarget\hypertarget{sourcepointer-1}{}\fi
\noindent\textit{Source reference.} Original lemma 01XB: sources/stacks-chow/tags.txt, source lines 2512--2512. Navigation pointer for \textbackslash{}ref\{coherent-lemma-quasi-coherent-affine-cohomology-zero\} used inside the counted span (line 1361); the argument that consumes it is retained in the span.\par\medskip
\ifdefined\hypertarget\hypertarget{sourcepointer-2}{}\fi
\noindent\textit{Source reference.} Original proposition 03BM: sources/stacks-chow/tags.txt, source lines 4353--4353. Navigation pointer for \textbackslash{}ref\{groupoids-proposition-finite-flat-equivalence\} used inside the counted span (line 1354); the argument that consumes it is retained in the span.\par\medskip
\ifdefined\hypertarget\hypertarget{sourcepointer-3}{}\fi
\noindent\textit{Source reference.} Original lemma 0728: sources/stacks-spaces-cohomology/pinned-original-source.tex, source lines 1217--1217. Navigation pointer for \textbackslash{}ref\{lemma-alternating-spectral-sequence\} used inside the counted span (line 1345); the argument that consumes it is retained in the span.\par\medskip
\ifdefined\hypertarget\hypertarget{sourcepointer-4}{}\fi
\noindent\textit{Source reference.} Original lemma 072A: sources/stacks-spaces-cohomology/pinned-original-source.tex, source lines 1305--1305. Navigation pointer for \textbackslash{}ref\{lemma-quasi-coherent-twist\} used inside the counted span (line 1363); the argument that consumes it is retained in the span.\par\medskip
\ifdefined\hypertarget\hypertarget{sourcepointer-5}{}\fi
\noindent\textit{Source reference.} Original section 071Y: sources/stacks-spaces-cohomology/pinned-original-source.tex, source lines 288--288. Navigation pointer for \textbackslash{}ref\{section-higher-direct-image\} used inside the counted span (line 1356); the argument that consumes it is retained in the span.\par\medskip
\ifdefined\hypertarget\hypertarget{sourcepointer-6}{}\fi
\noindent\textit{Source reference.} Original lemma 03KR: sources/stacks-chow/tags.txt, source lines 4687--4687. Navigation pointer for \textbackslash{}ref\{spaces-morphisms-lemma-compose-after-separated\} used inside the counted span (line 1348); the argument that consumes it is retained in the span.\par\medskip
\begin{proposition}
\label{proposition-vanishing}
Let $S$ be a scheme. Let $X$ be an algebraic space over $S$.
Assume $X$ is quasi-compact and separated.
Let $U$ be an affine scheme, and let
$f : U \to X$ be a surjective \'etale morphism.
Let $d$ be an upper bound for the size of the fibres of
$|U| \to |X|$. Then for any quasi-coherent $\mathcal{O}_X$-module $\mathcal{F}$
we have $H^q(X, \mathcal{F}) = 0$ for $q \geq d$.
\end{proposition}

\begin{proof}
We will use the spectral sequence of
Lemma \ref{lemma-alternating-spectral-sequence}.
The lemma applies since $f$ is separated as $U$ is separated, see
Morphisms of Spaces, Lemma
\ref{spaces-morphisms-lemma-compose-after-separated}.
Since $X$ is separated the scheme $U \times_X \ldots \times_X U$ is a closed
subscheme of
$U \times_{\Spec(\mathbf{Z})} \ldots \times_{\Spec(\mathbf{Z})} U$
hence is affine. Thus $W_p$ is affine. Hence $U_p = W_p/S_{p + 1}$ is an
affine scheme by
Groupoids, Proposition \ref{groupoids-proposition-finite-flat-equivalence}.
The discussion in
Section \ref{section-higher-direct-image}
shows that cohomology of quasi-coherent sheaves on $W_p$ (as an algebraic
space) agrees with the cohomology of the corresponding quasi-coherent
sheaf on the underlying affine scheme, hence vanishes in positive degrees by
Cohomology of Schemes, Lemma
\ref{coherent-lemma-quasi-coherent-affine-cohomology-zero}.
By
Lemma \ref{lemma-quasi-coherent-twist}
the sheaves
$\mathcal{F}|_{U_p} \otimes_\mathbf{Z} \underline{\mathbf{Z}}(\chi_p)$
are quasi-coherent. Hence
$H^q(W_p,
\mathcal{F}|_{U_p} \otimes_\mathbf{Z} \underline{\mathbf{Z}}(\chi_p))$
is zero when $q > 0$. By our definition of the integer $d$ we see that
$W_p = \emptyset$ for $p \geq d$. Hence also
$H^0(W_p,
\mathcal{F}|_{U_p} \otimes_\mathbf{Z} \underline{\mathbf{Z}}(\chi_p))$
is zero when $p \geq d$.
This proves the proposition.
\end{proof}

\end{document}
